\section{New explicit contraction schedules}
\label{sec:theory}

\subsection{Polynomial space: use the side-degree capacity at several depths}

The root pays the symmetric high-degree constraint
$Q_0^2\Delta\ll s$, while subsequent calls need only
$Q_d\Delta\ll s$ for their old side. Their fresh side is much smaller
by \eqref{eq:side-invariant}. The asymmetric packet already exploits this
distinction, and the round-two packet already changes the root alphabet.
Here the cumulative capacity is used at several depths, with progressively
smaller tail alphabets.

Choose the following rational retained-range coefficients:
\begin{equation}\label{eq:ps-schedule}
\begin{array}{c|ccccc}
 d&0&1&2&3&4+k\ (k\ge0)\\\hline
 r_d&13&14&8&6&4\\
 \eta_d&27/2000&51/4000&141/20000&107/20000&(27/6400)(5/8)^k.
\end{array}
\end{equation}
Use $u_d=\floor{\eta_d\lambda}$, the buffer \eqref{eq:buffer}, and leaf
threshold $H=\ceil{s/\lambda}$. All parameters are computable by integer
arithmetic. The cumulative leading ranges are
\begin{equation}\label{eq:ps-capacities}
 A_0=\frac{27}{2000},\quad A_1=\frac{21}{800},\quad
 A_2=\frac{333}{10000},\quad A_3=\frac{773}{20000},\quad
 A_\infty=\frac{499}{10000}.
\end{equation}

\begin{lemma}[Every packing gate has a strict margin]
\label{lem:ps-gates}
At the root, $2A_0\log_2r_0<1/10$. At all later depths,
$A_d\log_2r_d<1/10$. The three conditions in
\eqref{eq:packing-gates} hold uniformly at active logarithmic depths,
and all new tensor degrees stay below $\Delta$.
\end{lemma}
\begin{proof}
The required scalar inequalities have exact integer certificates
\[
 13^{27}<2^{100},\quad 14^{21}<2^{80},\quad
 999<1000,\quad 6^{773}<2^{2000},\quad 998<1000.
\]
They correspond in order to the root and the four tail alphabet regimes.
At the root they give $Q_0^2\Delta=s^{1-c_0+o(1)}$ for fixed $c_0>0$.
Below the root they give $Q_dR_d=s^{1-c_d+o(1)}$ with a uniform positive
minimum margin. The fresh-side exponent satisfies
\[
 2A_d\log_2r_d+A_{d-1}<\frac15+\frac{499}{10000}
 =\frac{2499}{10000}<1.
\]
The same bound controls $L_dQ_d^2$. The added buffers cost only
$O((\log\lambda)^2)=o(\lambda)$ at logarithmic depth, while active
calls have $s/M<\lambda$. These logarithmic factors are absorbed by the
strict power margins. Finally $2L_d=s^{0.0499+o(1)}<\Delta$.
\end{proof}

\paragraph{Noncircular depth argument.}
First prove the preceding estimates for every hypothetical active depth
\[
 d\le\ceil{\log_{5/3}\lambda}+C_0
\]
with a fixed sufficiently large $C_0$. In that range,
\[
 u_d\ge\lambda^{\epsilon_0+o(1)},\qquad
 \epsilon_0=1-\frac{\ln(8/5)}{\ln(5/3)}>0,
\]
where positivity follows from $24<25$. Thus floors are uniformly
negligible, all taper ranges diverge, and the packing and precision
estimates are valid throughout the hypothetical range. For sufficiently
large $s$, every child has at most $3M/5$ wires, because $r_d\ge4$ and
padding tends uniformly to one. No active path can then pass the stated
depth. This establishes the depth bound without assuming the desired
shrinkage before verifying packing.

Using the sharper accumulated padding estimate gives
\begin{equation}\label{eq:ps-size}
 M_d\le(1+o(1))s f_d,
 \quad (f_0,f_1,f_2,f_3)=\left(1,\frac2{13},\frac2{91},\frac1{182}\right),
 \quad f_{4+k}=\frac1{546}\,2^{-k}.
\end{equation}
Apply \eqref{eq:local-ps} and sum the logarithms of branching counts along
a path. The tail work ratio is
$(1/2)/\sqrt{5/8}=\sqrt{2/5}<1$. Therefore the exact leading coefficient is
\begin{align}
 C_{\rm ps}={}&\frac{K(13)}{\sqrt{27/2000}}
 +\frac{2K(14)}{13\sqrt{51/4000}}
 +\frac{2K(8)}{91\sqrt{141/20000}}\notag\\
 &+\frac{K(6)}{182\sqrt{107/20000}}
 +\frac{K(4)}{546\sqrt{27/6400}\,(1-\sqrt{2/5})}.
 \label{eq:ps-constant}
\end{align}
The exact rational interval checker proves
\begin{equation}\label{eq:ps-interval}
 63.85220339054<C_{\rm ps}<63.85220339056<64.
\end{equation}
It uses rational bounds from Machin's arctangent formula for $\pi$, a
positive series for $\ln2$, and integer-square-root brackets. A
floating-point parameter search suggested the schedule but is not an
input to the proof.

Polynomial choices of interpolation node, residue-field evaluation, and
child group over $O(\log\lambda)$ levels contribute
$O(\lambda\log\lambda)$ to the logarithm of total work. Leaf summation
costs $2^{O(s/\lambda)}$. Original entry oracles cost
$2^{O(s^{9/10})}$; new packet entries cost at most
$\exp(O(L_d\log s))=\exp(s^{0.0499+o(1)})$. All are lower order than
$\mathfrak E(s)$. Polynomial-bit descriptions, arithmetic and precision,
together with depth-first traversal, preserve polynomial space. The
description-length factor stays a fixed polynomial because branching
depends on $s$, not on $\mathcal L$. This proves
Theorem~\ref{thm:ps}.

\subsection{A stronger comparison than the displayed prior coefficients}

The incoming asymmetric schedule uses $r=16$, root range $\lambda/81$,
and a geometric tail of the same total range. Composing it with full-block
modes multiplies its slope-three coefficient by $\sqrt3/2$, giving
$66.6051434076\ldots$. Its fixed-alphabet capacity boundary and optimal
tail allocation give the stronger derived value
\begin{equation}\label{eq:prior-fixed}
 K(r)\sqrt{20\log_2r}
 \left[1+\frac{2/r}{(1-(2/r)^{2/3})^{3/2}}\right],
\end{equation}
which at $r=16$ is $66.1927236936\ldots$.

For a root alphabet $r_0$ and a constant tail alphabet $r_1$, set
\[
 x_0=\frac{.05}{\log_2r_0},\qquad
 U=\frac{.1}{\log_2r_1}-x_0>0.
\]
The corresponding optimally allocated geometric tail has coefficient
\[
 \frac{K(r_0)}{\sqrt{x_0}}+
 \frac{(2/r_0)K(r_1)}{\sqrt U\,(1-(2/r_1)^{2/3})^{3/2}}.
\]
At $(r_0,r_1)=(17,13)$ it is about $66.12047546$. These formulas give
feasible limiting comparisons by introducing a vanishing reserve in the
capacities; they do not assert global minima over all schedules. The new
strict rational coefficient in \eqref{eq:ps-interval} is below all three
comparisons. Its gain is about $4.13\%$ against the direct composed
certificate and $3.43\%$ against the stronger root/tail comparison.

\subsection{Cached contraction: pay for the group functions}
\label{sec:cache-theory}

The cached method is inherited from the joint-Cauchy packet \cite{joint}.
Partition a mode tuple by its exact total weight on every directed pair
of groups. With fixed $r$ and polynomial precision, there are only
polynomially many weight arrays, although the polynomial degree can be
large. For one array, let $h_e$ be the base-two logarithm of the number
of exact-weight tuples on pair $e$. Polynomial-size dynamic programming
over modes and weights supplies exact counts and rank/unrank operations;
an efficient hidden oracle is not assumed.

Put
\begin{equation}\label{eq:cache-K}
 \kappa=\frac{\pi\sqrt2}{\ln2},\qquad
 \mathcal K(M,u)=\kappa\frac{M}{\sqrt u}.
\end{equation}
From \eqref{eq:pair-families}, the slope-four generating function, and
the global weight budget, the common-tilt argument gives for fixed
nonnegative pair weights $z_e$,
\begin{equation}\label{eq:weighted-entropy}
 \sum_ez_e h_e\le
 \left(\sqrt{\frac{\sum_ez_e^2}{r^2}}+o(1)\right)\mathcal K(M,u).
\end{equation}
Using one common tilt matters: it includes zero and small individual
weight budgets without incorrectly applying a separate large-budget
asymptotic to each pair.

Cache each group's function of its incident pair indices, obtaining
each value by a genuine recursive child contraction. Partition the
groups into three blocks. Sum internal indices to form block boundary
tables, then contract the resulting triangle using division-free
rectangular Strassen multiplication. For block size $k$, the numbers
of incident and boundary directed pairs are respectively
\[
 I_r(k)=2rk-k^2-k,\qquad J_r(k)=2k(r-k).
\]
If the three inter-block classes have sizes $v_1,v_2,v_3$, the rectangular
matrix exponent follows from
\[
 H_1+H_2+H_3-(3-\omega)\min(H_1,H_2,H_3),\qquad \omega=\log_2 7.
\]
Writing this as the maximum of three weighted sums and using
\eqref{eq:weighted-entropy} gives the matrix coefficient
\begin{equation}\label{eq:matrix-b}
 b_r=\sqrt{\frac{v_1+v_2+v_3-[1-(\omega-2)^2]\min_jv_j}{r^2}}.
\end{equation}
The cache, block-construction, and memory coefficients are
\begin{equation}\label{eq:cache-factors}
 a_r=\frac{\sqrt{2(r-1)}}r,\qquad
 B_r=\max\left\{b_r,\frac{\sqrt{\max_k I_r(k)}}r\right\},\qquad
 j_r=\frac{\sqrt{\max_k J_r(k)}}r.
\end{equation}
Here $B_r$ must include both block-table construction and matrix work.
Omitting the former would be a mistake for the root schedule below.

Up to the already controlled polynomial and lower-order factors, the
correct recurrence is
\begin{equation}\label{eq:cache-recurrence}
 T_d\le2^{a_{r_d}\mathcal K_d}\,r_dT_{d+1}
           +2^{B_{r_d}\mathcal K_d}.
\end{equation}
Thus caching does not grant free group tensors. The first term pays for
all recursive evaluations, and the second pays for the tables and their
contraction. The inherited sequential-record implementation, reblocking
and sorting show that polynomial overhead suffices in a bit-cost model;
an exponential random-access table is not treated as a unit-cost oracle.
Strassen uses only additions, subtractions and multiplications, so it is
valid in the source's residue ring with nilpotents.

\subsection{Why the stronger fixed-alphabet baseline must be included}

At the asymmetric packet's displayed $r=16$ schedule, the new per-pair
counts and full-block cached method give time coefficient
$44.5375718550\ldots$ and space coefficient $39.4959495337\ldots$.
But one can retune the existing fixed-alphabet family. At $r=8$, root
capacity $x_0=1/60$, and a tail of the same total capacity, the appropriate
recurrence is root dominated for a range of geometric tail allocations.
Its limiting time coefficient is
\begin{equation}\label{eq:r8-baseline}
 C_{8,\partial}=\kappa\sqrt{60}
  \sqrt{\frac{30+12(\omega-2)^2}{64}}
  =38.1677456970\ldots,
\end{equation}
with space coefficient $33.9927010637\ldots$. Vanishing capacity reserve
handles the equality boundary. This is the more demanding comparator
used here, rather than only the $44.54$ displayed-schedule composition.

\subsection{The new strict root-7/tail-9 schedule}

Use seven groups at the root, partitioned as $(2,2,3)$, and nine groups
at every later depth, partitioned as $(3,3,3)$. Set
\begin{equation}\label{eq:new-cache-schedule}
 \eta_0=\frac{89}{5000},\qquad
 \eta_d=\frac{1233}{125000}\left(\frac7{25}\right)^{d-1}\quad(d\ge1),
 \qquad u_d=\floor{\eta_d\lambda}.
\end{equation}
Keep the buffer \eqref{eq:buffer} and $H=\ceil{s/\lambda}$.
The tail total is $137/10000$, and the total cumulative capacity is
\[
 A_\infty=\frac{89}{5000}+\frac{137}{10000}
 =\frac{63}{2000}=0.0315.
\]
The exact inequalities
\begin{equation}\label{eq:cache-gates}
 7^{89}<2^{250},\qquad 9^{63}<2^{200}
\end{equation}
give the strict root gate $2\eta_0\log_2 7<1/10$ and tail gate
$A_\infty\log_2 9<1/10$. The fresh-side and padding exponents are
less than $1/5+A_\infty=0.2315<1$. New tensor degrees are at most
$s^{0.0315+o(1)}<\Delta$.

Verify these bounds first through hypothetical depth
$\ceil{\log_4\lambda}+O(1)$. Tail ranges satisfy
\[
 u_d\ge\lambda^{1-\ln(25/7)/\ln4+o(1)}\longrightarrow\infty
\]
because $25/7<4$. Eventually the root child shrinkage is at most $1/3$
and every tail child shrinkage at most $1/4$, while the sharp factors
are $2/7$ and $2/9$ respectively. This closes the same noncircular
depth bootstrap as before and validates the uniform precision estimates.
Accumulated padding errors give
\[
 M_0\le s,\qquad
 M_d\le(1+o(1))s\frac27\left(\frac29\right)^{d-1}\quad(d\ge1).
\]

Write $\xi=\omega-2=\log_2(7/4)$. At the root,
\begin{equation}\label{eq:root-cache-factors}
 a_7=\frac{\sqrt{12}}7,\quad
 b_7=\frac{\sqrt{24+8\xi^2}}7,\quad
 B_7=\frac{\sqrt{30}}7,\quad j_7=\frac{\sqrt{24}}7.
\end{equation}
Here $b_7<B_7$: root block construction costs more than root matrix
multiplication. In the tail,
\begin{equation}\label{eq:tail-cache-factors}
 a_9=\frac49,\quad
 B_9=b_9=\frac{\sqrt{4+2\xi^2}}3,\quad j_9=\frac23.
\end{equation}
The tail block-construction coefficient $\sqrt{42}/9$ is smaller than
$B_9$. The tail work ratio is
\begin{equation}\label{eq:cache-gamma}
 \gamma=\frac{2/9}{\sqrt{7/25}}=\frac{10}{9\sqrt7}<1.
\end{equation}
Moreover $a_9/(1-\gamma)=0.766231\ldots<B_9=0.767655\ldots$.
This strict inequality ensures that all deeper cache evaluations are
covered by the first tail node's local table bound.

Let
\[
 K_0=\frac{\kappa}{\sqrt{89/5000}},\qquad
 K_1=\frac{(2/7)\kappa}{\sqrt{1233/125000}}.
\]
Unrolling \eqref{eq:cache-recurrence}, the normalized total time
coefficient is at most
\begin{equation}\label{eq:cache-final}
 C_{\rm cache}=\max\{B_7K_0,\ a_7K_0+B_9K_1\}
 =a_7K_0+B_9K_1.
\end{equation}
The root table term is $37.59170256\ldots$; the second term, including
recursive child work, is $37.9301090829\ldots$ and is the bottleneck.
The memory coefficient is
\begin{equation}\label{eq:cache-memory}
 C_{\rm mem}=\max\{j_7K_0,j_9K_1\}=j_7K_0
 =33.6230409217\ldots.
\end{equation}
Ancestor and descendant table sizes add, so their leading space exponents
take a maximum, rather than a sum. There are only logarithmically many
stack levels, and the tail $\mathcal K_d$ decreases geometrically.

The supplied checker proves these comparisons with rational intervals,
including a positive margin in $a_9/(1-\gamma)<B_9$ and the complete
root maximum. It encloses
\begin{align*}
 37.93010908288&<C_{\rm cache}<37.93010908289<38,\\
 33.62304092165&<C_{\rm mem}<33.62304092167<34.
\end{align*}
The same lower-order costs already checked for Theorem~\ref{thm:ps}
are harmless here, with an even smaller packet exponent. This proves
Theorem~\ref{thm:cached}. The time coefficient improves by about $0.62\%$
over \eqref{eq:r8-baseline}; the comparison is a strict improvement of
certificates, not a prediction of a $0.62\%$ runtime reduction.

\subsection{Consequences for SAT and the unresolved order improvement}

Let $n$ be the declared variable count, $m$ the number of clauses and
$k$ a fixed width. The unchanged reduction and the density refinements
give the polynomial-space envelope
\begin{equation}\label{eq:sat-envelope}
 (n+m)^{O_k(1)}\,
 2^{O_k\left(\frac n{\sqrt{\log n}}
      \min\{\log(2+m/n),\log\log n\}\right)}.
\end{equation}
Small sizes are understood with the usual harmless logarithm offsets.
At fixed density this is $2^{O_{k,D}(n/\sqrt{\log n})}$ up to polynomial
factors. At unrestricted density it has the
$n\log\log n/\sqrt{\log n}$ numerator/denominator profile. The new
engine preserves these orders and improves its specified wire-level
certificate. Converting it to a numerical coefficient per variable
requires the reduction's actual wire and width constants. In particular,
writing $2^{64n/\sqrt{\log n}}$ as a universal SAT consequence would
be unsupported.

The cached version can use the same reductions at its different space
cost. Direct clause-network bounds and prior low-degree specializations
remain options in a portfolio. The separate manuscript's unspecified
constants cannot be ranked numerically from the supplied information.

The persistent obstacle to a smaller logarithmic order is unchanged:
the mode weights still grow linearly, their weighted partitions have
the $M/\sqrt u$ entropy scale, and packing keeps $u=O(\log s)$.
The existing scoped barriers do not prove a lower bound for SAT; they
show why the present retuning and exact arithmetic transformations do
not establish $n/\log n$ or eliminate the dense reduction's extra
$\log\log n$. A further order improvement still requires an actual
summation/factorization interface or a different representation,
with its oracle and wire costs paid explicitly.
